\chapter{Games}\label{ch:iv1:games}

\begin{ivtrains}
Thinking through what the other side will do. Take-away games yield to a losing-position table and to nim-sums, sequential games to backward induction, and simultaneous games to the rule that a mix must leave the opponent indifferent. Toy poker games and auctions put these together.
\end{ivtrains}

\begin{question}[$\star$ \iqroles{QR, QT, Trader, Analyst} \iqfirm{any} \iqlevel{junior}]\label{iq:iv1:games:1}
There is a pile of 97 stones. Two players take turns removing between 1 and 6 stones, and whoever takes the last stone wins. Should you go first, and what is your first move? How do things change if whoever takes the last stone loses?
\end{question}

\begin{question}[$\star$ \iqroles{QR, QT, Trader} \iqfirm{any} \iqlevel{junior}]\label{iq:iv1:games:2}
There are two piles, of 5 and 8 stones. Two players take turns, and each turn removes any positive number of stones from one pile. Whoever takes the last stone wins. Who wins, and how? What if the piles are equal to begin with?
\end{question}

\begin{question}[$\star$ \iqroles{QR, QT, Trader, Analyst} \iqfirm{any} \iqlevel{junior}]\label{iq:iv1:games:3}
You and I each show a coin at the same moment. If both show heads, I pay you \$3. If both show tails, I pay you \$1. If they differ, you pay me \$2. What should each of us do, and what is this game worth to you per round?
\end{question}

\begin{question}[$\star$ \iqroles{QR, QT, Trader, Analyst} \iqfirm{any} \iqlevel{junior}]\label{iq:iv1:games:4}
Everyone in a room writes a whole number from 0 to 100. The winner is the person closest to two-thirds of the average of all the numbers. If everybody is fully rational and knows that everybody else is too, what is written? What would you write in a real room, and why?
\end{question}

\begin{question}[$\star\star$ \iqroles{QR, QT, Trader} \iqfirm{any} \iqlevel{junior}]\label{iq:iv1:games:5}
In a game of rock, paper, scissors, a win with rock pays \$1, a win with paper pays \$2, and a win with scissors pays \$3, and a tie pays nothing. What is the equilibrium mix? Which option is played most often, and why does that look odd at first?
\end{question}

\begin{question}[$\star\star$ \iqroles{QR, QT, Trader, QD} \iqfirm{any} \iqlevel{junior}]\label{iq:iv1:games:6}
There are three piles, of 3, 5 and 7 stones. Each turn removes any positive number of stones from one pile, and whoever takes the last stone wins. Should you go first? List all winning first moves, and prove the rule that you use.
\end{question}

\begin{question}[$\star\star$ \iqroles{QR, QT, QD} \iqfirm{any} \iqlevel{junior}]\label{iq:iv1:games:7}
In one pile, a player may remove 1, 3 or 4 stones each turn, and whoever takes the last stone wins. Compute the Grundy values for piles of up to 8 stones. Now the game is played on two such piles, of 10 and 6, and each turn is a move in one pile. Who wins, and with what move?
\end{question}

\begin{question}[$\star\star$ \iqroles{QR, QT, Trader, Analyst} \iqfirm{any} \iqlevel{junior}]\label{iq:iv1:games:8}
Five pirates, ranked A to E, share 100 gold coins. The highest-ranked living pirate proposes a split, and all living pirates vote. If at least half agree, the split stands; otherwise, the proposer is thrown overboard, and the next one proposes. Each cares first about survival, then about coins, and then prefers to see someone thrown overboard. What does A propose? How does that change if a proposal needs more than half of the votes?
\end{question}

\begin{question}[$\star\star$ \iqroles{QR, QT, Trader} \iqfirm{market maker, hedge fund} \iqlevel{junior}]\label{iq:iv1:games:9}
Two traders divide \$1 by alternating offers: A offers a split, B accepts or rejects; if B rejects, B offers, then A accepts or rejects, and so on. Each round of delay multiplies the value of the money by $0.9$. If there are at most three rounds, and after a rejection in the third both get nothing, what does A get? And if there is no limit on rounds?
\end{question}

\begin{question}[$\star\star$ \iqroles{QR, QT, Trader} \iqfirm{any} \iqlevel{junior}]\label{iq:iv1:games:10}
A penalty taker aims left or right, and a keeper dives left or right, at the same moment. The chance of a goal when the taker aims left is $0.6$ if the keeper dives left and $0.95$ if the keeper dives right. When the taker aims right it is $0.9$ if the keeper dives left and $0.5$ if the keeper dives right. Find the equilibrium and the chance of a goal. The taker then improves, and the $0.9$ becomes $0.95$. Does the taker aim right more or less often?
\end{question}

\begin{question}[$\star\star$ \iqroles{QR, QT, Trader, Analyst} \iqfirm{market maker, crypto} \iqlevel{junior}]\label{iq:iv1:games:11}
$n$ bidders have private values that are independent and uniform on $[0, 1]$. In a second-price auction, what should each bid, and why? In a first-price auction, find the symmetric equilibrium bid. What revenue does the seller expect in each?
\end{question}

\begin{question}[$\star\star$ \iqroles{QR, QT, Trader, Analyst} \iqfirm{hedge fund, market maker} \iqlevel{junior}]\label{iq:iv1:games:12}
A task is worth \$4 to each of $n$ people, and it gets done if at least one of them does it. Doing it costs the doer \$1. They all decide at the same moment, with no talk. In the symmetric equilibrium, where each does it with the same probability, what is that probability for $n = 2$ and for $n = 3$? How likely is it that nobody does it, and what happens as $n$ grows?
\end{question}

\begin{question}[$\star\star$ \iqroles{QR, QT} \iqfirm{any} \iqlevel{junior}]\label{iq:iv1:games:13}
A chocolate bar is a grid of $m$ by $n$ squares. A move picks a remaining square and eats it and every square above it and to its right. The player who eats the bottom-left square loses. Show that the first player has a winning strategy whenever the bar has more than one square, without finding what the winning move is. Give the explicit strategy on a square bar.
\end{question}

\begin{question}[$\star\star\star$ \iqroles{QR, QT, Trader} \iqfirm{hedge fund, market maker} \iqlevel{junior, experienced}]\label{iq:iv1:games:14}
A crew of pirates, ranked from top to bottom, share 100 gold coins. The highest-ranked living pirate proposes a split, and all living pirates vote, the proposer included. If at least half agree, the split stands; otherwise, the proposer is thrown overboard, and the next in rank proposes. Each pirate cares first about survival, then about coins, and then prefers to see someone thrown overboard. With up to 200 pirates, the first proposer survives. What happens to the first proposer with 201, 202, 203 and 204 pirates? With 300 pirates, how many are thrown overboard?
\end{question}

\begin{question}[$\star\star\star$ \iqroles{QR, QT, Trader} \iqfirm{market maker, hedge fund} \iqlevel{junior, experienced}]\label{iq:iv1:games:15}
A three-card game is played with an ace, a king and a queen. X and Y each ante \$1 and receive one card. X may check, and then the higher card wins the \$2 pot, or bet \$1. Facing a bet, Y may fold, and then X takes the pot, or call, and then the higher card wins the pot of \$4. Find the equilibrium: how often does X bluff with the queen, and how often does Y call with the king? What is the game worth to X?
\end{question}

\begin{question}[$\star\star\star$ \iqroles{QR, QT, Trader} \iqfirm{any} \iqlevel{junior, experienced}]\label{iq:iv1:games:16}
Two bidders have independent private values, uniform on $[0, 1]$, and the seller's own value is 0. In a second-price auction with a reserve price $r$, nobody can win below $r$. Which reserve maximises the seller's expected revenue, and what is the revenue? Does the best reserve change with more bidders? What does a reserve do for the total surplus?
\end{question}

\begin{question}[$\star\star\star$ \iqroles{QR, QT, QD} \iqfirm{any} \iqlevel{junior, experienced}]\label{iq:iv1:games:17}
Coins sit on the steps of a staircase, with 3, 1, 4, 1 and 5 coins on steps 1 to 5. A move takes any positive number of coins from one step and puts them on the step below; from step 1 they leave the staircase. Whoever makes the last move wins. Who wins, and what is the first move? Explain why half of the steps can be ignored.
\end{question}

\begin{question}[$\star\star\star$ \iqroles{QR, QT, QD} \iqfirm{any} \iqlevel{junior, experienced}]\label{iq:iv1:games:18}
There are two piles of stones. A move removes any positive number from one pile, or the same positive number from both piles. Whoever takes the last stone wins. List the losing positions up to piles of 15. Is the position with piles 9 and 15 a win for the player to move, and with piles 10 and 15? What is the closed form for the losing positions?
\end{question}

\begin{question}[$\star\star\star\star$ \iqroles{QR, QT} \iqfirm{market maker, hedge fund} \iqlevel{experienced}]\label{iq:iv1:games:19}
A poker game uses three cards, a jack, a queen and a king. Each player antes \$1 and gets one card. Player 1 checks or bets \$1. After a check, Player 2 checks, for a showdown, or bets \$1, and then Player 1 folds or calls. After a bet, Player 2 folds or calls. Find the value of the game to Player 1, and describe a set of equilibrium strategies. Which hands bluff, and which player has the edge?
\end{question}

\begin{question}[$\star\star\star\star$ \iqroles{QR, QT} \iqfirm{hedge fund, market maker} \iqlevel{experienced}]\label{iq:iv1:games:20}
X and Y each ante \$1 and each receive a private hand uniform on $[0, 1]$, the higher hand being better. X may check, and then the higher hand takes the \$2 pot, or bet \$1. Facing a bet, Y folds, and X takes the pot, or calls, and the higher hand takes the \$4 pot. Find the equilibrium, in the form: X bets with very weak hands and very strong hands, and Y calls with a strong enough hand. What is the game worth to X?
\end{question}
